% Generated draft: requires mathematical and source review.
\documentclass[11pt]{article}
\usepackage[margin=25mm]{geometry}
\usepackage{fontspec}
\setmainfont{Latin Modern Roman}
\usepackage{amsmath,amssymb,mathtools,unicode-math}
\setmathfont{Latin Modern Math}
\usepackage{enumitem,needspace,xurl,hyperref}
\setlength{\parindent}{0pt}
\setlength{\parskip}{5pt}
\setlength{\emergencystretch}{3em}
\newcommand{\N}{\mathbb{N}}
\newcommand{\Z}{\mathbb{Z}}
\newcommand{\Q}{\mathbb{Q}}
\newcommand{\R}{\mathbb{R}}
\newcommand{\C}{\mathbb{C}}
\newcommand{\F}{\mathbb{F}}
\newcommand{\A}{\mathbb{A}}
\newcommand{\PP}{\mathbb{P}}
\newcommand{\E}{\mathbb{E}}
\newcommand{\eps}{\varepsilon}
\newcommand{\dd}{\,\mathrm{d}}
\begin{document}
\section*{UnsolvedMath ID: 11000182}
\textbf{Source problem code:} catalogue entry
\subsection{Definitions and mathematical statement}
Retain the source's symplectic notation \(\omega\), \(K(\omega)\), and \(\mathcal C_{\chi,\theta}\). In the usual closed symplectic \(4\)-manifold setting, for a compatible almost-complex structure \(J\) one sets \(K(\omega)=-c_1(TM,J)\in H^2(M;\Z)\), and
\[
 [\omega]\cdot K(\omega)=\langle[\omega]\smile K(\omega),[M]\rangle.
\]
The supplied excerpt does not explicitly give the dimension, compactness, minimality, canonical-class convention, or definitions of \(\chi,\theta,\mathcal C_{\chi,\theta}\); this standard interpretation is conditional on its missing setup.

\emph{Problem 4.} Reproduce the classification of manifolds satisfying
\[
 [\omega]\cdot K(\omega)<0
\]
by studying the sets \(\mathcal C_{\chi,\theta}\) for \(\theta<0\). Both the classification result to be recovered and the sets providing the requested method must be specified from the preceding source discussion. The problem is not merely to state a classification by any unrelated method, and the excerpt does not identify \(\theta\) with a particular characteristic number.

The source then raises a network of questions in the borderline case \(K(\omega)\cdot[\omega]=0\), and in the stronger case \(K(\omega)=0\). It does not identify vanishing of the pairing with vanishing of the class. Integral vanishing, vanishing in real cohomology, and torsion of a canonical class are also distinct and must be separated. It reports that the only examples then known with the latter condition were standard complex tori, certain other torus bundles over tori, and K3 surfaces. This is a dated source assessment, not an exhaustive classification theorem or a statement about every torus bundle. The specific subsequent question introduced by ``So we have'' is not included.
\subsection{Short English statement}
Recover the negative-canonical-pairing manifold classification using the source's sets \(\mathcal C_{\chi,\theta}\) with \(\theta<0\). The following zero-pairing and zero-canonical-class discussion is separate context, with only a historical list of known examples.
\subsection{Sources}
\begin{itemize}
\item Local UnsolvedMath registry, permanent ID 11000182, source code catalogue entry, Problem 4, PDF page 178. No individual proposer or separate classification citation is supplied.
\item Farb, \emph{Problems on Mapping Class Groups and Related Topics}. Supplied URL: \url{https://www.math.uchicago.edu/~farb/papers/mcgbook.pdf}; recorded access 2026-07-29; extraction pdf-font-structure. The volume contains 329 explicitly typeset headings across chapters, with repeated numbering disambiguated by page.
\item The source itself lists standard complex tori, certain other torus bundles over tori, and K3 surfaces as its then-known zero-canonical-class examples, without a separate reference or the following promised question.
\item Supplied status open and OPEN-TRIAGE. The research summary says only ``Cannot verify.'' Current primary-literature status and redistribution terms for the publicly accessible author/publisher copy remain NEEDS\_REVIEW.
\end{itemize}
\end{document}
